Se considera el problema
\[
y''-y'=f(t),\qquad y(0)=0,\qquad t\ge0,
\]
en variables adimensionales. La pendiente inicial puede elegirse. Se estudian dos entradas por separado: \(f(t)=\cos t\) y \(f(t)=1\).

Elija y justifique una vía de resolución antes de efectuar los cálculos. Para cada entrada, determine si existe una pendiente inicial que mantenga \(y\) acotada durante todo \(t\ge0\). Si existe, obtenga esa pendiente y la solución correspondiente; si no, explique por qué ningún valor de la pendiente puede conseguirlo. Su argumento debe distinguir los términos que se pueden modificar mediante las condiciones iniciales de los impuestos por la entrada. Aquí «acotada» significa que existe una constante finita \(B\) tal que \(|y(t)|\le B\) para todo \(t\ge0\).
